\section{Exercices d'application}

\rubrique{Repère, coordonnées, distance et milieu}

\exo{}\diff{1}
Dans un repère orthonormé $(O,\vec{i},\vec{j})$, on donne les points $A(2;3)$, $B(-1;4)$ et $C(5;-2)$.
\begin{enumerate}[label=\alph*)]
\item Calculer les distances $AB$, $AC$ et $BC$.
\item Le triangle $ABC$ est-il rectangle ?
\item Déterminer les coordonnées du milieu $I$ de $[BC]$.
\end{enumerate}

\exo{}\diff{1}
On considère les points $E(-3;2)$ et $F(5;-4)$.
\begin{enumerate}[label=\alph*)]
\item Calculer les coordonnées du point $M$ tel que $\overrightarrow{EM} = \frac{1}{2}\overrightarrow{EF}$.
\item Déterminer les coordonnées du point $N$ symétrique de $E$ par rapport à $F$.
\end{enumerate}

\exo{}\diff{1}
Dans un repère orthonormé, on donne $A(1;2)$, $B(4;3)$ et $C(3;-1)$.
\begin{enumerate}[label=\alph*)]
\item Montrer que $ABC$ est un triangle isocèle.
\item Calculer les coordonnées du centre du cercle circonscrit au triangle $ABC$.
\end{enumerate}

\rubrique{Vecteurs et colinéarité}

\exo{}\diff{1}
Soient $\vec{u}(2; -3)$ et $\vec{v}(-4; 6)$.
\begin{enumerate}[label=\alph*)]
\item Les vecteurs $\vec{u}$ et $\vec{v}$ sont-ils colinéaires ? Justifier.
\item Déterminer le réel $k$ tel que $\vec{v}=k\vec{u}$.
\end{enumerate}

\exo{}\diff{2}
On donne $A(-2;1)$, $B(3;4)$, $C(5;-1)$ et $D(0;-4)$.
\begin{enumerate}[label=\alph*)]
\item Montrer que $ABCD$ est un parallélogramme.
\item Calculer les coordonnées du point $E$ tel que $ABCE$ soit un parallélogramme.
\end{enumerate}

\exo{}\diff{2}
Dans un repère, on considère les points $P(1;2)$, $Q(4;5)$ et $R(7;8)$.
\begin{enumerate}[label=\alph*)]
\item Montrer que $P$, $Q$ et $R$ sont alignés.
\item Déterminer le réel $t$ tel que $\overrightarrow{PR}=t\overrightarrow{PQ}$.
\end{enumerate}

\rubrique{Équations de droites}

\exo{}\diff{1}
Dans un repère orthonormé, déterminer une équation cartésienne de la droite $(d)$ passant par $A(2;-1)$ et de vecteur directeur $\vec{v}(3;4)$.

\exo{}\diff{1}
Soit $(d)$ la droite d'équation $2x-3y+6=0$.
\begin{enumerate}[label=\alph*)]
\item Donner un vecteur directeur de $(d)$.
\item Déterminer l'équation réduite de $(d)$.
\item Le point $B(5;4)$ appartient-il à $(d)$ ?
\end{enumerate}

\exo{}\diff{2}
On donne les droites $(d_1): y=2x-1$ et $(d_2): -x+2y-3=0$.
\begin{enumerate}[label=\alph*)]
\item Les droites $(d_1)$ et $(d_2)$ sont-elles parallèles ? Justifier.
\item Déterminer les coordonnées de leur point d'intersection $I$.
\end{enumerate}

\exo{}\diff{2}
Dans un repère orthonormé, on considère les points $A(1;2)$, $B(3;5)$ et $C(4;1)$.
\begin{enumerate}[label=\alph*)]
\item Déterminer une équation de la droite $(AB)$.
\item Déterminer une équation de la droite passant par $C$ et parallèle à $(AB)$.
\item Déterminer une équation de la droite passant par $C$ et perpendiculaire à $(AB)$.
\end{enumerate}

\rubrique{Équation d'un cercle}

\exo{}\diff{1}
Déterminer l'équation du cercle de centre $\Omega(2;-3)$ et de rayon $R=5$.

\exo{}\diff{2}
On considère l'équation $x^2+y^2-4x+6y-12=0$.
\begin{enumerate}[label=\alph*)]
\item Montrer que cette équation est celle d'un cercle dont on précisera le centre et le rayon.
\item Le point $A(5;1)$ appartient-il à ce cercle ?
\end{enumerate}

\exo{}\diff{2}
Déterminer l'équation du cercle passant par les points $A(1;2)$, $B(3;4)$ et $C(5;2)$.

\rubrique{Positions relatives droite-cercle et distance point-droite}

\exo{}\diff{2}
On donne le cercle $(\mathcal{C})$ d'équation $(x-1)^2+(y+2)^2=25$ et la droite $(d): 3x-4y+9=0$.
\begin{enumerate}[label=\alph*)]
\item Calculer la distance du centre de $(\mathcal{C})$ à la droite $(d)$.
\item En déduire la position relative de $(d)$ et $(\mathcal{C})$.
\end{enumerate}

\exo{}\diff{3}
Soit $(\mathcal{C})$ le cercle d'équation $x^2+y^2-2x+4y-20=0$ et $(d)$ la droite passant par $A(1;3)$ et de vecteur directeur $\vec{v}(2;1)$.
\begin{enumerate}[label=\alph*)]
\item Déterminer une équation cartésienne de $(d)$.
\item Étudier la position relative de $(d)$ et $(\mathcal{C})$.
\item Déterminer les coordonnées des points d'intersection éventuels.
\end{enumerate}

\exo{}\diff{3}
Dans un repère orthonormé, on considère les points $A(2;1)$, $B(5;4)$ et $C(3;6)$.
\begin{enumerate}[label=\alph*)]
\item Déterminer une équation du cercle $(\mathcal{C})$ de diamètre $[AB]$.
\item Montrer que $C$ appartient à $(\mathcal{C})$.
\item Déterminer une équation de la tangente à $(\mathcal{C})$ au point $C$.
\end{enumerate}

\section{Exercices d'entraînement}

\rubrique{Repère, coordonnées, distance et milieu}

\exo{}\diff{1}
Dans un repère orthonormé, on donne $A(-1;3)$ et $B(4;-2)$. Calculer $AB$ et les coordonnées du milieu de $[AB]$.

\exo{}\diff{1}
Soient $M(2;5)$ et $N(-3;1)$. Déterminer les coordonnées du point $P$ tel que $M$ soit le milieu de $[NP]$.

\exo{}\diff{2}
On considère les points $A(1;1)$, $B(4;2)$ et $C(3;5)$. Calculer les distances $AB$, $AC$ et $BC$. Quelle est la nature du triangle $ABC$ ?

\exo{}\diff{2}
Dans un repère orthonormé, on donne $A(-2;3)$, $B(4;1)$ et $C(2;-3)$. Montrer que $ABC$ est un triangle rectangle.

\rubrique{Vecteurs et colinéarité}

\exo{}\diff{1}
Soient $\vec{u}(3; -2)$ et $\vec{v}(-6; 4)$. Les vecteurs $\vec{u}$ et $\vec{v}$ sont-ils colinéaires ?

\exo{}\diff{2}
On donne $A(1;2)$, $B(3;5)$, $C(4;1)$ et $D(2;-2)$. Montrer que $ABCD$ est un parallélogramme.

\exo{}\diff{2}
Soient $E(-1;2)$, $F(3;4)$ et $G(7;6)$. Montrer que $E$, $F$ et $G$ sont alignés.

\exo{}\diff{2}
Dans un repère, on considère les points $P(2;1)$, $Q(5;3)$ et $R(8;5)$. Déterminer le réel $k$ tel que $\overrightarrow{PR}=k\overrightarrow{PQ}$.

\rubrique{Équations de droites}

\exo{}\diff{1}
Déterminer une équation cartésienne de la droite passant par $A(1;-2)$ et de vecteur directeur $\vec{v}(-2;3)$.

\exo{}\diff{1}
Soit $(d): 3x+2y-5=0$. Donner un vecteur directeur de $(d)$ et son équation réduite.

\exo{}\diff{2}
On donne $(d_1): y=-2x+3$ et $(d_2): 4x+2y-7=0$. Les droites $(d_1)$ et $(d_2)$ sont-elles parallèles ?

\exo{}\diff{2}
Déterminer l'intersection des droites $(d_1): x-2y+3=0$ et $(d_2): 2x+y-4=0$.

\exo{}\diff{2}
Dans un repère orthonormé, on donne $A(2;1)$, $B(4;3)$ et $C(1;4)$. Déterminer une équation de la droite passant par $C$ et parallèle à $(AB)$.

\exo{}\diff{3}
Soient $A(1;2)$, $B(3;5)$ et $C(4;1)$. Déterminer une équation de la hauteur issue de $A$ dans le triangle $ABC$.

\rubrique{Équation d'un cercle}

\exo{}\diff{1}
Déterminer l'équation du cercle de centre $\Omega(-1;4)$ et de rayon $R=3$.

\exo{}\diff{2}
Montrer que $x^2+y^2+6x-8y+9=0$ est l'équation d'un cercle. Préciser son centre et son rayon.

\exo{}\diff{2}
Déterminer l'équation du cercle de diamètre $[AB]$ avec $A(1;3)$ et $B(5;-1)$.

\exo{}\diff{3}
Déterminer l'équation du cercle passant par $A(2;1)$, $B(4;